\cotefile{87}
\scanpage{29}
\[
C\subset E \qquad
\begin{aligned}
&C\ \text{partie polyédrale de l'espace affine }E\\
&\qquad \dim C=d
\end{aligned}
\]

$A_C=\mathrm{Aff}(C)$ espace affine engendré par $C$

$C^0=\mathrm{Int}_{A_C}(C)$ est la facette maximale (= de dim maximale), $C=\overline{C^0}$.

Les autres facettes \inserted{fermées} sont les intersections \inserted{$H\cap C$ avec $C$} d'hyperplans $H$
de $A_C$ ne rencontrant pas $C^0$ \deleted{\unsure{avec le}}. Les facettes
\deleted{fermées \unsure{telles q|telles que}} de dim $d-1$ correspondent \inserted{$1-1$} aux
$H$ tels que \deleted{\unsure{l'}}l'intérieur de $H\cap C$ dans $H$ soit $\ne\varnothing$,
et les intérieurs en question sont les facettes \inserted{\unsure{ouv}} de
dim $d-1$. Soit $C^1$ leur réunion.

Les facettes fermées de dim $\leq d-2$ sont les int. avec
$C$ de sous-espaces affines de dim $d-2$ de $A_C$ ne
rencontrant pas $C^0\cup C^1\ldots$

Les facettes fermées de dim $\leq d-i$ sont les int.
\inserted{$F\cap C$} avec $C$ des sous-espaces \inserted{\unsure{aff}} $F$ de $A_C$ de dim
$d-i$ \inserted{(i.e. \unsure{codim|codim.} $i$)} qui ne rencontrent pas $C^0\cup C^1\cup\ldots\cup C^{i-1}$.
Les facettes fermées de dim $d-i$ sont en
corr. $1-1$ avec les $F$ tels que $\mathrm{int}_{F}\,F\cap C\ne\varnothing$,
et \inserted{les} $\mathrm{int}_F\,F\cap C$ sont les facettes de dim $d-i$.

Intersection de deux facettes fermées est une
facette fermée [c'est une partie convexe de $C$ et
c'est une réunion de facettes, il reste \unsure{à prouver|: prouver}):

% Vertical pen stroke to the left of the next two manuscript lines.
\uline{Prop} Les facettes \inserted{fermées} de $C$ sont les réunions
de facette qui sont convexes et fermées.

Soit
% archive sheet number: 28
